% Options for packages loaded elsewhere
% Options for packages loaded elsewhere
\PassOptionsToPackage{unicode}{hyperref}
\PassOptionsToPackage{hyphens}{url}
\PassOptionsToPackage{dvipsnames,svgnames,x11names}{xcolor}
%
\documentclass[
  letterpaper,
  DIV=11,
  numbers=noendperiod]{scrartcl}
\usepackage{xcolor}
\usepackage{amsmath,amssymb}
\setcounter{secnumdepth}{-\maxdimen} % remove section numbering
\usepackage{iftex}
\ifPDFTeX
  \usepackage[T1]{fontenc}
  \usepackage[utf8]{inputenc}
  \usepackage{textcomp} % provide euro and other symbols
\else % if luatex or xetex
  \usepackage{unicode-math} % this also loads fontspec
  \defaultfontfeatures{Scale=MatchLowercase}
  \defaultfontfeatures[\rmfamily]{Ligatures=TeX,Scale=1}
\fi
\usepackage{lmodern}
\ifPDFTeX\else
  % xetex/luatex font selection
\fi
% Use upquote if available, for straight quotes in verbatim environments
\IfFileExists{upquote.sty}{\usepackage{upquote}}{}
\IfFileExists{microtype.sty}{% use microtype if available
  \usepackage[]{microtype}
  \UseMicrotypeSet[protrusion]{basicmath} % disable protrusion for tt fonts
}{}
\makeatletter
\@ifundefined{KOMAClassName}{% if non-KOMA class
  \IfFileExists{parskip.sty}{%
    \usepackage{parskip}
  }{% else
    \setlength{\parindent}{0pt}
    \setlength{\parskip}{6pt plus 2pt minus 1pt}}
}{% if KOMA class
  \KOMAoptions{parskip=half}}
\makeatother
% Make \paragraph and \subparagraph free-standing
\makeatletter
\ifx\paragraph\undefined\else
  \let\oldparagraph\paragraph
  \renewcommand{\paragraph}{
    \@ifstar
      \xxxParagraphStar
      \xxxParagraphNoStar
  }
  \newcommand{\xxxParagraphStar}[1]{\oldparagraph*{#1}\mbox{}}
  \newcommand{\xxxParagraphNoStar}[1]{\oldparagraph{#1}\mbox{}}
\fi
\ifx\subparagraph\undefined\else
  \let\oldsubparagraph\subparagraph
  \renewcommand{\subparagraph}{
    \@ifstar
      \xxxSubParagraphStar
      \xxxSubParagraphNoStar
  }
  \newcommand{\xxxSubParagraphStar}[1]{\oldsubparagraph*{#1}\mbox{}}
  \newcommand{\xxxSubParagraphNoStar}[1]{\oldsubparagraph{#1}\mbox{}}
\fi
\makeatother

\usepackage{color}
\usepackage{fancyvrb}
\newcommand{\VerbBar}{|}
\newcommand{\VERB}{\Verb[commandchars=\\\{\}]}
\DefineVerbatimEnvironment{Highlighting}{Verbatim}{commandchars=\\\{\}}
% Add ',fontsize=\small' for more characters per line
\usepackage{framed}
\definecolor{shadecolor}{RGB}{241,243,245}
\newenvironment{Shaded}{\begin{snugshade}}{\end{snugshade}}
\newcommand{\AlertTok}[1]{\textcolor[rgb]{0.68,0.00,0.00}{#1}}
\newcommand{\AnnotationTok}[1]{\textcolor[rgb]{0.37,0.37,0.37}{#1}}
\newcommand{\AttributeTok}[1]{\textcolor[rgb]{0.40,0.45,0.13}{#1}}
\newcommand{\BaseNTok}[1]{\textcolor[rgb]{0.68,0.00,0.00}{#1}}
\newcommand{\BuiltInTok}[1]{\textcolor[rgb]{0.00,0.23,0.31}{#1}}
\newcommand{\CharTok}[1]{\textcolor[rgb]{0.13,0.47,0.30}{#1}}
\newcommand{\CommentTok}[1]{\textcolor[rgb]{0.37,0.37,0.37}{#1}}
\newcommand{\CommentVarTok}[1]{\textcolor[rgb]{0.37,0.37,0.37}{\textit{#1}}}
\newcommand{\ConstantTok}[1]{\textcolor[rgb]{0.56,0.35,0.01}{#1}}
\newcommand{\ControlFlowTok}[1]{\textcolor[rgb]{0.00,0.23,0.31}{\textbf{#1}}}
\newcommand{\DataTypeTok}[1]{\textcolor[rgb]{0.68,0.00,0.00}{#1}}
\newcommand{\DecValTok}[1]{\textcolor[rgb]{0.68,0.00,0.00}{#1}}
\newcommand{\DocumentationTok}[1]{\textcolor[rgb]{0.37,0.37,0.37}{\textit{#1}}}
\newcommand{\ErrorTok}[1]{\textcolor[rgb]{0.68,0.00,0.00}{#1}}
\newcommand{\ExtensionTok}[1]{\textcolor[rgb]{0.00,0.23,0.31}{#1}}
\newcommand{\FloatTok}[1]{\textcolor[rgb]{0.68,0.00,0.00}{#1}}
\newcommand{\FunctionTok}[1]{\textcolor[rgb]{0.28,0.35,0.67}{#1}}
\newcommand{\ImportTok}[1]{\textcolor[rgb]{0.00,0.46,0.62}{#1}}
\newcommand{\InformationTok}[1]{\textcolor[rgb]{0.37,0.37,0.37}{#1}}
\newcommand{\KeywordTok}[1]{\textcolor[rgb]{0.00,0.23,0.31}{\textbf{#1}}}
\newcommand{\NormalTok}[1]{\textcolor[rgb]{0.00,0.23,0.31}{#1}}
\newcommand{\OperatorTok}[1]{\textcolor[rgb]{0.37,0.37,0.37}{#1}}
\newcommand{\OtherTok}[1]{\textcolor[rgb]{0.00,0.23,0.31}{#1}}
\newcommand{\PreprocessorTok}[1]{\textcolor[rgb]{0.68,0.00,0.00}{#1}}
\newcommand{\RegionMarkerTok}[1]{\textcolor[rgb]{0.00,0.23,0.31}{#1}}
\newcommand{\SpecialCharTok}[1]{\textcolor[rgb]{0.37,0.37,0.37}{#1}}
\newcommand{\SpecialStringTok}[1]{\textcolor[rgb]{0.13,0.47,0.30}{#1}}
\newcommand{\StringTok}[1]{\textcolor[rgb]{0.13,0.47,0.30}{#1}}
\newcommand{\VariableTok}[1]{\textcolor[rgb]{0.07,0.07,0.07}{#1}}
\newcommand{\VerbatimStringTok}[1]{\textcolor[rgb]{0.13,0.47,0.30}{#1}}
\newcommand{\WarningTok}[1]{\textcolor[rgb]{0.37,0.37,0.37}{\textit{#1}}}

\usepackage{longtable,booktabs,array}
\newcounter{none} % for unnumbered tables
\usepackage{calc} % for calculating minipage widths
% Correct order of tables after \paragraph or \subparagraph
\usepackage{etoolbox}
\makeatletter
\patchcmd\longtable{\par}{\if@noskipsec\mbox{}\fi\par}{}{}
\makeatother
% Allow footnotes in longtable head/foot
\IfFileExists{footnotehyper.sty}{\usepackage{footnotehyper}}{\usepackage{footnote}}
\makesavenoteenv{longtable}
\usepackage{graphicx}
\makeatletter
\newsavebox\pandoc@box
\newcommand*\pandocbounded[1]{% scales image to fit in text height/width
  \sbox\pandoc@box{#1}%
  \Gscale@div\@tempa{\textheight}{\dimexpr\ht\pandoc@box+\dp\pandoc@box\relax}%
  \Gscale@div\@tempb{\linewidth}{\wd\pandoc@box}%
  \ifdim\@tempb\p@<\@tempa\p@\let\@tempa\@tempb\fi% select the smaller of both
  \ifdim\@tempa\p@<\p@\scalebox{\@tempa}{\usebox\pandoc@box}%
  \else\usebox{\pandoc@box}%
  \fi%
}
% Set default figure placement to htbp
\def\fps@figure{htbp}
\makeatother





\setlength{\emergencystretch}{3em} % prevent overfull lines

\providecommand{\tightlist}{%
  \setlength{\itemsep}{0pt}\setlength{\parskip}{0pt}}



 


\KOMAoption{captions}{tableheading}
\makeatletter
\@ifpackageloaded{caption}{}{\usepackage{caption}}
\AtBeginDocument{%
\ifdefined\contentsname
  \renewcommand*\contentsname{Table of contents}
\else
  \newcommand\contentsname{Table of contents}
\fi
\ifdefined\listfigurename
  \renewcommand*\listfigurename{List of Figures}
\else
  \newcommand\listfigurename{List of Figures}
\fi
\ifdefined\listtablename
  \renewcommand*\listtablename{List of Tables}
\else
  \newcommand\listtablename{List of Tables}
\fi
\ifdefined\figurename
  \renewcommand*\figurename{Figure}
\else
  \newcommand\figurename{Figure}
\fi
\ifdefined\tablename
  \renewcommand*\tablename{Table}
\else
  \newcommand\tablename{Table}
\fi
}
\@ifpackageloaded{float}{}{\usepackage{float}}
\floatstyle{ruled}
\@ifundefined{c@chapter}{\newfloat{codelisting}{h}{lop}}{\newfloat{codelisting}{h}{lop}[chapter]}
\floatname{codelisting}{Listing}
\newcommand*\listoflistings{\listof{codelisting}{List of Listings}}
\makeatother
\makeatletter
\makeatother
\makeatletter
\@ifpackageloaded{caption}{}{\usepackage{caption}}
\@ifpackageloaded{subcaption}{}{\usepackage{subcaption}}
\makeatother
\usepackage{bookmark}
\IfFileExists{xurl.sty}{\usepackage{xurl}}{} % add URL line breaks if available
\urlstyle{same}
% fallback for those not using the hyperref driver hyperxmp:
\makeatletter
\@ifundefined{xmpquote}{\newcommand{\xmpquote}[1]{#1}}{}
\makeatother
\hypersetup{
  pdftitle={Simulations \& Tidyverse lab: Intermediate Version},
  pdfauthor={Lumi Kang},
  colorlinks=true,
  linkcolor={blue},
  filecolor={Maroon},
  citecolor={Blue},
  urlcolor={Blue},
  pdfcreator={LaTeX via pandoc}}


\title{Simulations \& Tidyverse lab: Intermediate Version}
\usepackage{etoolbox}
\makeatletter
\providecommand{\subtitle}[1]{% add subtitle to \maketitle
  \apptocmd{\@title}{\par {\large #1 \par}}{}{}
}
\makeatother
\subtitle{PSYC 201A, Data Simulation Lab 1}
\author{Lumi Kang}
\date{2026-10-07}
\begin{document}
\maketitle

\renewcommand*\contentsname{Table of contents}
{
\setcounter{tocdepth}{2}
\tableofcontents
}

\begin{quote}
\textbf{Goals:} (1) Understand how normal, binomial, and lognormal
distributions describe behavioral data (accuracy, RTs, individual
variability). (2) Practice generating and visualizing synthetic datasets
using tidyverse functions. (3) Learn how to write reproducible lab
reports with inline statistics in Quarto.
\end{quote}

\begin{quote}
\textbf{Workflow:} This intermediate version provides starter code and
scaffolding to help you complete each task. Fill in the missing pieces
and adapt the provided code.
\end{quote}

\subsection{Setup}\label{setup}

If you finished getting-started-with-r, the tidyverse is already
installed. If not, run this chunk once, by hand (green arrow at its top
right). It is set to \texttt{eval:\ false} so it does not run every time
you render.

\begin{Shaded}
\begin{Highlighting}[]
\FunctionTok{install.packages}\NormalTok{(}\StringTok{"tidyverse"}\NormalTok{)}
\end{Highlighting}
\end{Shaded}

\texttt{set.seed()} makes the random numbers come out the same every
time you render, so the numbers you write about in your text do not
change under you. Change the seed and every simulated value below will
change.

\begin{Shaded}
\begin{Highlighting}[]
\FunctionTok{library}\NormalTok{(tidyverse)}
\FunctionTok{set.seed}\NormalTok{(}\DecValTok{2025}\NormalTok{)}
\end{Highlighting}
\end{Shaded}

\subsection{A. Normal distribution (worked
example)}\label{a.-normal-distribution-worked-example}

Here, let's try to simulate some values from a normal distribution with
a mean of 0 and a SD of 1.

The \texttt{rnorm} function is used below here.

\begin{Shaded}
\begin{Highlighting}[]
\CommentTok{\# you can play around with these}
\NormalTok{n }\OtherTok{\textless{}{-}} \DecValTok{300} \CommentTok{\# number of datapoints}
\NormalTok{mu }\OtherTok{\textless{}{-}} \DecValTok{0} \CommentTok{\# mean of the distribution}
\NormalTok{sd\_norm }\OtherTok{\textless{}{-}} \DecValTok{1} \CommentTok{\# sd of the distribution}

\CommentTok{\# rnorm is your simulation function here}
\CommentTok{\# take a look at how it\textquotesingle{}s used by typing \textasciigrave{}?rnorm\textasciigrave{} in the Console}
\DocumentationTok{\#\# note that this function is taking the inputs you specified above {-}{-} it doesn\textquotesingle{}t matter what these are called}
\NormalTok{norm\_df }\OtherTok{\textless{}{-}} \FunctionTok{tibble}\NormalTok{(}\AttributeTok{sim\_values =} \FunctionTok{rnorm}\NormalTok{(n, mu, sd\_norm))}

\CommentTok{\# so you know, this code would do the same thing}
\CommentTok{\# norm\_df \textless{}{-} tibble(sim\_values = rnorm(300, 0, 1))}

\CommentTok{\# okay, check out the top of this dataframe}
\FunctionTok{head}\NormalTok{(norm\_df)}
\end{Highlighting}
\end{Shaded}

\begin{verbatim}
# A tibble: 6 x 1
  sim_values
       <dbl>
1     0.621 
2     0.0356
3     0.773 
4     1.27  
5     0.371 
6    -0.163 
\end{verbatim}

Calculate some basic summaries of this dataframe and print them out

\begin{Shaded}
\begin{Highlighting}[]
\CommentTok{\# }\AlertTok{TODO}\CommentTok{: Calculate mean, sd, and n for norm\_df}
\CommentTok{\# Hint: Use summarise() with mean(), sd(), and n()}
\CommentTok{\# Start with this structure and fill in the missing parts:}

\NormalTok{norm\_df\_summary }\OtherTok{\textless{}{-}}\NormalTok{ norm\_df }\SpecialCharTok{|\textgreater{}} 
  \FunctionTok{summarise}\NormalTok{(}
    \AttributeTok{mean =} \FunctionTok{mean}\NormalTok{(sim\_values),  }\CommentTok{\# Fill in the column name}
    \AttributeTok{sd =} \FunctionTok{sd}\NormalTok{(sim\_values),      }\CommentTok{\# Fill in the column name  }
    \AttributeTok{n =} \FunctionTok{n}\NormalTok{()            }\CommentTok{\# This one is already done!}
\NormalTok{  )}

\CommentTok{\# Print the results}
\FunctionTok{print}\NormalTok{(norm\_df\_summary)}
\end{Highlighting}
\end{Shaded}

\begin{verbatim}
# A tibble: 1 x 3
     mean    sd     n
    <dbl> <dbl> <int>
1 0.00433 0.987   300
\end{verbatim}

\begin{Shaded}
\begin{Highlighting}[]
\CommentTok{\# What happens if you increase the "n"? How do these summaries change?}
\CommentTok{\# It seems like the mean becomes negative and gets smaller.}
\end{Highlighting}
\end{Shaded}

I'm going to now plot a histogram of these values. We're going to spend
a lot of time on figuring out how to plot things later -- but here you
can start learning the basics.

The first argument to the ggplot function is your dataframe -- here,
\texttt{norm\_df}. The \texttt{aes} is where you specify what is on the
x-axis, y-axis, colors, etc. These values come from the columns in
\texttt{norm\_df}.

\texttt{Alpha} is transparency. \texttt{binwidth} is how big the
histogram bins are.

\begin{Shaded}
\begin{Highlighting}[]
\FunctionTok{ggplot}\NormalTok{(norm\_df, }\FunctionTok{aes}\NormalTok{(}\AttributeTok{x=}\NormalTok{sim\_values)) }\SpecialCharTok{+}
  \FunctionTok{geom\_histogram}\NormalTok{(}\AttributeTok{binwidth =} \FloatTok{0.25}\NormalTok{, }\AttributeTok{alpha =} \FloatTok{0.85}\NormalTok{) }\SpecialCharTok{+}
  \FunctionTok{labs}\NormalTok{(}\AttributeTok{title =} \StringTok{"Simulated values"}\NormalTok{, }\AttributeTok{x =} \StringTok{"Value"}\NormalTok{, }\AttributeTok{y =} \StringTok{"Count"}\NormalTok{)}
\end{Highlighting}
\end{Shaded}

\pandocbounded{\includegraphics[keepaspectratio]{distributions-lab-intermediate_files/figure-pdf/unnamed-chunk-4-1.pdf}}

\textbf{Exercise:} change \texttt{mu} and \texttt{sd\_norm} to two other
combos (e.g., 10 \& 2; 10 \& 5).

\begin{Shaded}
\begin{Highlighting}[]
\NormalTok{mu2 }\OtherTok{=} \DecValTok{10}
\NormalTok{sd\_norm2 }\OtherTok{=} \DecValTok{2}
\NormalTok{norm\_df2 }\OtherTok{\textless{}{-}} \FunctionTok{tibble}\NormalTok{(}\AttributeTok{sim\_values =} \FunctionTok{rnorm}\NormalTok{(n, mu2, sd\_norm2))}

\NormalTok{norm\_df2\_summary }\OtherTok{\textless{}{-}}\NormalTok{ norm\_df2 }\SpecialCharTok{|\textgreater{}} 
  \FunctionTok{summarise}\NormalTok{(}
    \AttributeTok{mean =} \FunctionTok{mean}\NormalTok{(sim\_values),  }\CommentTok{\# Fill in the column name}
    \AttributeTok{sd =} \FunctionTok{sd}\NormalTok{(sim\_values),      }\CommentTok{\# Fill in the column name  }
    \AttributeTok{n =} \FunctionTok{n}\NormalTok{()            }\CommentTok{\# This one is already done!}
\NormalTok{  )}

\NormalTok{mu3 }\OtherTok{=} \DecValTok{10}
\NormalTok{sd\_norm3 }\OtherTok{=} \DecValTok{5}
\NormalTok{norm\_df3 }\OtherTok{\textless{}{-}} \FunctionTok{tibble}\NormalTok{(}\AttributeTok{sim\_values =} \FunctionTok{rnorm}\NormalTok{(n, mu3, sd\_norm3))}
\NormalTok{norm\_df3\_summary }\OtherTok{\textless{}{-}}\NormalTok{ norm\_df3 }\SpecialCharTok{|\textgreater{}} 
  \FunctionTok{summarise}\NormalTok{(}
    \AttributeTok{mean =} \FunctionTok{mean}\NormalTok{(sim\_values),  }\CommentTok{\# Fill in the column name}
    \AttributeTok{sd =} \FunctionTok{sd}\NormalTok{(sim\_values),      }\CommentTok{\# Fill in the column name }
    \AttributeTok{n =} \FunctionTok{n}\NormalTok{()            }\CommentTok{\# This one is already done!}
\NormalTok{  )}
\end{Highlighting}
\end{Shaded}

Write one sentence with an inline mean (e.g., 0).

\begin{Shaded}
\begin{Highlighting}[]
\FunctionTok{round}\NormalTok{(}\FunctionTok{mean}\NormalTok{(norm\_df2}\SpecialCharTok{$}\NormalTok{sim\_values), }\DecValTok{2}\NormalTok{)}
\end{Highlighting}
\end{Shaded}

\begin{verbatim}
[1] 10.02
\end{verbatim}

\begin{Shaded}
\begin{Highlighting}[]
\FunctionTok{round}\NormalTok{(}\FunctionTok{mean}\NormalTok{(norm\_df3}\SpecialCharTok{$}\NormalTok{sim\_values), }\DecValTok{2}\NormalTok{)}
\end{Highlighting}
\end{Shaded}

\begin{verbatim}
[1] 10.02
\end{verbatim}

\textbf{Bonus:} Compute another kind of summary statistic on the
dataset.

\begin{Shaded}
\begin{Highlighting}[]
\NormalTok{norm\_df2\_summary }\OtherTok{\textless{}{-}}\NormalTok{ norm\_df2 }\SpecialCharTok{|\textgreater{}} 
  \FunctionTok{summarise}\NormalTok{(}
    \AttributeTok{mean =} \FunctionTok{mean}\NormalTok{(sim\_values),  }\CommentTok{\# Fill in the column name}
    \AttributeTok{sd =} \FunctionTok{sd}\NormalTok{(sim\_values),      }\CommentTok{\# Fill in the column name }
    \AttributeTok{n =} \FunctionTok{n}\NormalTok{(),            }\CommentTok{\# This one is already done!}
    \AttributeTok{median =} \FunctionTok{median}\NormalTok{(sim\_values),}
    \AttributeTok{min =} \FunctionTok{min}\NormalTok{(sim\_values),}
    \AttributeTok{max =} \FunctionTok{max}\NormalTok{(sim\_values)}
\NormalTok{  )}

\NormalTok{norm\_df3\_summary }\OtherTok{\textless{}{-}}\NormalTok{ norm\_df3 }\SpecialCharTok{|\textgreater{}}
  \FunctionTok{summarise}\NormalTok{(}
    \AttributeTok{mean =} \FunctionTok{mean}\NormalTok{(sim\_values),  }\CommentTok{\# Fill in the column name}
    \AttributeTok{sd =} \FunctionTok{sd}\NormalTok{(sim\_values),      }\CommentTok{\# Fill in the column name }
    \AttributeTok{n =} \FunctionTok{n}\NormalTok{(),            }\CommentTok{\# This one is already done!}
    \AttributeTok{median =} \FunctionTok{median}\NormalTok{(sim\_values),}
    \AttributeTok{min =} \FunctionTok{min}\NormalTok{(sim\_values),}
    \AttributeTok{max =} \FunctionTok{max}\NormalTok{(sim\_values)}
\NormalTok{  )}
\end{Highlighting}
\end{Shaded}

\textbf{Bonus:} Simulate another distribution and also plot it on the
same histogram (you can make it a different color)

\begin{Shaded}
\begin{Highlighting}[]
\FunctionTok{ggplot}\NormalTok{(norm\_df2, }\FunctionTok{aes}\NormalTok{(}\AttributeTok{x=}\NormalTok{sim\_values)) }\SpecialCharTok{+}
  \FunctionTok{geom\_histogram}\NormalTok{(}\AttributeTok{binwidth =} \FloatTok{0.5}\NormalTok{, }\AttributeTok{alpha =} \FloatTok{0.5}\NormalTok{, }\AttributeTok{fill =} \StringTok{"steelblue"}\NormalTok{) }\SpecialCharTok{+}
  \FunctionTok{labs}\NormalTok{(}\AttributeTok{title =} \StringTok{"Simulated values"}\NormalTok{, }\AttributeTok{x =} \StringTok{"Value"}\NormalTok{, }\AttributeTok{y =} \StringTok{"Count"}\NormalTok{)}
\end{Highlighting}
\end{Shaded}

\pandocbounded{\includegraphics[keepaspectratio]{distributions-lab-intermediate_files/figure-pdf/unnamed-chunk-8-1.pdf}}

\begin{Shaded}
\begin{Highlighting}[]
\FunctionTok{ggplot}\NormalTok{(norm\_df3, }\FunctionTok{aes}\NormalTok{(}\AttributeTok{x=}\NormalTok{sim\_values)) }\SpecialCharTok{+}
  \FunctionTok{geom\_histogram}\NormalTok{(}\AttributeTok{binwidth =} \FloatTok{0.5}\NormalTok{, }\AttributeTok{alpha =} \FloatTok{0.5}\NormalTok{, }\AttributeTok{fill =} \StringTok{"forestgreen"}\NormalTok{) }\SpecialCharTok{+}
  \FunctionTok{labs}\NormalTok{(}\AttributeTok{title =} \StringTok{"Simulated values"}\NormalTok{, }\AttributeTok{x =} \StringTok{"Value"}\NormalTok{, }\AttributeTok{y =} \StringTok{"Count"}\NormalTok{)}
\end{Highlighting}
\end{Shaded}

\pandocbounded{\includegraphics[keepaspectratio]{distributions-lab-intermediate_files/figure-pdf/unnamed-chunk-8-2.pdf}}

\subsection{B. Binomial data
(accuracy)}\label{b.-binomial-data-accuracy}

Okay, normally distributed data is great, but often not what we're
working with in practice. Often, we have accuracies or reaction times.

First, we'll think about accuracy. To do so, we'll practice
\texttt{rbinom()} and how to summarize it using some tidyverse
functions. You should be able to adapt the code we just worked through
but with a new function --

\subsubsection{B1 --- Bernoulli accuracy}\label{b1-bernoulli-accuracy}

Tasks:

\begin{itemize}
\tightlist
\item
  Simulate \texttt{n\ =\ 200} 0/1 outcomes with \texttt{p\ =\ 0.92} into
  \texttt{acc\_df}.
\item
  Compute \textbf{proportion correct} and a \textbf{count table} with
  \texttt{count()}.
\item
  Plot overall proportion with \texttt{stat\_summary()}.
\end{itemize}

\begin{Shaded}
\begin{Highlighting}[]
\CommentTok{\# }\AlertTok{TODO}\CommentTok{: create a simulated dataframe using rbinom (remember, you can look at the help for the function)}
\CommentTok{\# Hint: Use tibble() and rbinom(n=200, size=1, prob=0.92)}
\CommentTok{\# Complete this line:}
\NormalTok{acc\_df }\OtherTok{\textless{}{-}} \FunctionTok{tibble}\NormalTok{(}\AttributeTok{acc =} \FunctionTok{rbinom}\NormalTok{(}\DecValTok{200}\NormalTok{, }\AttributeTok{size =} \DecValTok{1}\NormalTok{, }\AttributeTok{prob =} \FloatTok{0.92}\NormalTok{))}


\CommentTok{\# }\AlertTok{TODO}\CommentTok{: summarise and count}
\CommentTok{\# Hint: Use summarise() to calculate proportion correct and n}
\CommentTok{\# Complete this structure:}
\NormalTok{acc\_df\_summary }\OtherTok{\textless{}{-}}\NormalTok{ acc\_df }\SpecialCharTok{|\textgreater{}}
  \FunctionTok{summarise}\NormalTok{(}
    \AttributeTok{prop\_correct =} \FunctionTok{mean}\NormalTok{(acc),  }\CommentTok{\# Fill in the column name}
    \AttributeTok{n =} \FunctionTok{n}\NormalTok{()                    }\CommentTok{\# This one is already done!}
\NormalTok{  )}

\CommentTok{\# Print the results}
\FunctionTok{print}\NormalTok{(acc\_df\_summary)}
\end{Highlighting}
\end{Shaded}

\begin{verbatim}
# A tibble: 1 x 2
  prop_correct     n
         <dbl> <int>
1          0.9   200
\end{verbatim}

\begin{Shaded}
\begin{Highlighting}[]
\CommentTok{\# Count table: how many 0s and 1s?}
\NormalTok{acc\_df }\SpecialCharTok{|\textgreater{}} \FunctionTok{count}\NormalTok{(acc)}
\end{Highlighting}
\end{Shaded}

\begin{verbatim}
# A tibble: 2 x 2
    acc     n
  <int> <int>
1     0    20
2     1   180
\end{verbatim}

\begin{Shaded}
\begin{Highlighting}[]
\DocumentationTok{\#\# TO DO: Make a histogram of these values}
\CommentTok{\# Hint: Use ggplot() with geom\_histogram()}
\CommentTok{\# Complete this structure:}
\FunctionTok{ggplot}\NormalTok{(acc\_df, }\FunctionTok{aes}\NormalTok{(}\AttributeTok{x =}\NormalTok{ acc)) }\SpecialCharTok{+}  \CommentTok{\# Fill in dataframe and column name}
  \FunctionTok{geom\_histogram}\NormalTok{(}\AttributeTok{binwidth =} \FloatTok{0.5}\NormalTok{, }\AttributeTok{alpha =} \FloatTok{0.85}\NormalTok{) }\SpecialCharTok{+}  \CommentTok{\# Choose appropriate values}
  \FunctionTok{labs}\NormalTok{(}\AttributeTok{title =} \StringTok{"Simulated accuracy values"}\NormalTok{, }\AttributeTok{x =} \StringTok{"Accuracy"}\NormalTok{, }\AttributeTok{y =} \StringTok{"Count"}\NormalTok{)}
\end{Highlighting}
\end{Shaded}

\pandocbounded{\includegraphics[keepaspectratio]{distributions-lab-intermediate_files/figure-pdf/unnamed-chunk-9-1.pdf}}

\begin{Shaded}
\begin{Highlighting}[]
\CommentTok{\# }\AlertTok{TODO}\CommentTok{: plot overall proportion }
\CommentTok{\# Hint:  use \textasciigrave{}stat\_summary(fun = mean)\textasciigrave{} for the mean}

\CommentTok{\# }\AlertTok{TODO}\CommentTok{: Plot some measure of variability (SD or 95\% CI)}
\CommentTok{\# Hint: stat\_summary(fun.data = mean\_se, fun.args = list(mult = 1.96), geom = "errorbar", width = 0.1) +}
\CommentTok{\# mean\_se(mult = 1.96) gives the mean +/{-} 1.96 standard errors, an approximate 95\% CI}
\CommentTok{\# You have to give stat\_summary an x{-}value, but since there isn\textquotesingle{}t really one here you can just give it a label (e.g., "overall")}

\CommentTok{\# Complete this structure:}
\FunctionTok{ggplot}\NormalTok{(acc\_df, }\FunctionTok{aes}\NormalTok{(}\AttributeTok{x =} \StringTok{\textquotesingle{}Overall\textquotesingle{}}\NormalTok{, }\AttributeTok{y =}\NormalTok{ acc)) }\SpecialCharTok{+}  \CommentTok{\# Fill in the column name}
  \FunctionTok{stat\_summary}\NormalTok{(}\AttributeTok{fun =}\NormalTok{ mean, }\AttributeTok{geom =} \StringTok{"point"}\NormalTok{, }\AttributeTok{size =} \DecValTok{3}\NormalTok{) }\SpecialCharTok{+}
  \FunctionTok{stat\_summary}\NormalTok{(}\AttributeTok{fun.data =}\NormalTok{ mean\_se, }\AttributeTok{fun.args =} \FunctionTok{list}\NormalTok{(}\AttributeTok{mult =} \FloatTok{1.96}\NormalTok{), }\AttributeTok{geom =} \StringTok{"errorbar"}\NormalTok{, }\AttributeTok{width =} \FloatTok{0.1}\NormalTok{) }\SpecialCharTok{+}
  \FunctionTok{labs}\NormalTok{(}\AttributeTok{x =} \ConstantTok{NULL}\NormalTok{, }\AttributeTok{y =} \StringTok{"Proportion correct"}\NormalTok{, }\AttributeTok{title =} \StringTok{"Accuracy"}\NormalTok{)}
\end{Highlighting}
\end{Shaded}

\pandocbounded{\includegraphics[keepaspectratio]{distributions-lab-intermediate_files/figure-pdf/unnamed-chunk-10-1.pdf}}

\subsubsection{B2 --- Two conditions + blocks
(tidy)}\label{b2-two-conditions-blocks-tidy}

Tasks:

\begin{itemize}
\tightlist
\item
  Build a tibble with \texttt{condition} in \{easy, hard\}, 100 trials
  each.
\item
  Use different probabilities (e.g., 0.97 vs 0.87).
\item
  Summarise \texttt{prop\_correct} by \texttt{condition\ ×\ block} and
  \texttt{pivot\_wider()}.
\item
  Bonus: Add a \texttt{participant} factor
\end{itemize}

\begin{Shaded}
\begin{Highlighting}[]
\CommentTok{\# }\AlertTok{TODO}\CommentTok{: simulate}
\CommentTok{\# Hint: Create a tibble with condition, block, and acc columns}
\CommentTok{\# Use rep() to repeat conditions and blocks appropriately}
\CommentTok{\# Use rbinom() with different probabilities for easy vs hard conditions}

\CommentTok{\# Complete this structure:}
\NormalTok{acc2 }\OtherTok{\textless{}{-}} \FunctionTok{tibble}\NormalTok{(}
  \AttributeTok{condition =} \FunctionTok{rep}\NormalTok{(}\FunctionTok{c}\NormalTok{(}\StringTok{"easy"}\NormalTok{, }\StringTok{"hard"}\NormalTok{), }\AttributeTok{each =} \DecValTok{100}\NormalTok{),  }\CommentTok{\# Fill in number of trials per condition}
  \AttributeTok{block =} \FunctionTok{rep}\NormalTok{(}\FunctionTok{rep}\NormalTok{(}\DecValTok{1}\SpecialCharTok{:}\DecValTok{4}\NormalTok{, }\AttributeTok{each =} \DecValTok{25}\NormalTok{), }\AttributeTok{times =} \DecValTok{2}\NormalTok{),   }\CommentTok{\# Fill in trials per block}
  \AttributeTok{acc =} \FunctionTok{c}\NormalTok{(}
    \FunctionTok{rbinom}\NormalTok{(}\DecValTok{100}\NormalTok{, }\DecValTok{1}\NormalTok{, }\FloatTok{0.97}\NormalTok{),  }\CommentTok{\# Fill in number of trials for easy condition}
    \FunctionTok{rbinom}\NormalTok{(}\DecValTok{100}\NormalTok{, }\DecValTok{1}\NormalTok{, }\FloatTok{0.87}\NormalTok{)   }\CommentTok{\# Fill in number of trials for hard condition}
\NormalTok{  )}
\NormalTok{)}

\CommentTok{\# Check per{-}block counts}
\NormalTok{acc2 }\SpecialCharTok{|\textgreater{}}
  \FunctionTok{add\_count}\NormalTok{(block, }\AttributeTok{name =} \StringTok{"n\_block"}\NormalTok{) }\SpecialCharTok{|\textgreater{}}
  \FunctionTok{distinct}\NormalTok{(block, n\_block)}
\end{Highlighting}
\end{Shaded}

\begin{verbatim}
# A tibble: 4 x 2
  block n_block
  <int>   <int>
1     1      50
2     2      50
3     3      50
4     4      50
\end{verbatim}

\begin{Shaded}
\begin{Highlighting}[]
\CommentTok{\# }\AlertTok{TODO}\CommentTok{: Summaries by condition × block}
\CommentTok{\# Hint: Use group\_by() and summarise()}
\CommentTok{\# Complete this structure:}
\NormalTok{acc\_by\_cb }\OtherTok{\textless{}{-}}\NormalTok{ acc2 }\SpecialCharTok{|\textgreater{}}
  \FunctionTok{group\_by}\NormalTok{(condition, block) }\SpecialCharTok{|\textgreater{}}  \CommentTok{\# Fill in grouping variables}
  \FunctionTok{summarise}\NormalTok{(}
    \AttributeTok{prop\_correct =} \FunctionTok{mean}\NormalTok{(acc),  }\CommentTok{\# Fill in column name}
    \AttributeTok{trials =} \FunctionTok{n}\NormalTok{(), }
    \AttributeTok{.groups =} \StringTok{"drop"}
\NormalTok{  )}

\CommentTok{\# }\AlertTok{TODO}\CommentTok{: Pivot wider}
\CommentTok{\# Hint: Use pivot\_wider() to make conditions into separate columns}
\CommentTok{\# Complete this structure:}
\NormalTok{acc\_wide }\OtherTok{\textless{}{-}}\NormalTok{ acc\_by\_cb }\SpecialCharTok{|\textgreater{}}
  \FunctionTok{pivot\_wider}\NormalTok{(}
    \AttributeTok{names\_from =}\NormalTok{ condition,     }\CommentTok{\# Fill in the column to become column names}
    \AttributeTok{values\_from =}\NormalTok{ prop\_correct     }\CommentTok{\# Fill in the column to become values}
\NormalTok{  )}

\CommentTok{\# View the results}
\FunctionTok{print}\NormalTok{(}\StringTok{"Summary by condition and block:"}\NormalTok{)}
\end{Highlighting}
\end{Shaded}

\begin{verbatim}
[1] "Summary by condition and block:"
\end{verbatim}

\begin{Shaded}
\begin{Highlighting}[]
\FunctionTok{print}\NormalTok{(acc\_by\_cb)}
\end{Highlighting}
\end{Shaded}

\begin{verbatim}
# A tibble: 8 x 4
  condition block prop_correct trials
  <chr>     <int>        <dbl>  <int>
1 easy          1         1        25
2 easy          2         0.92     25
3 easy          3         1        25
4 easy          4         0.96     25
5 hard          1         0.8      25
6 hard          2         0.96     25
7 hard          3         0.88     25
8 hard          4         0.8      25
\end{verbatim}

\begin{Shaded}
\begin{Highlighting}[]
\FunctionTok{print}\NormalTok{(}\StringTok{"Wide format:"}\NormalTok{)}
\end{Highlighting}
\end{Shaded}

\begin{verbatim}
[1] "Wide format:"
\end{verbatim}

\begin{Shaded}
\begin{Highlighting}[]
\FunctionTok{print}\NormalTok{(acc\_wide)}
\end{Highlighting}
\end{Shaded}

\begin{verbatim}
# A tibble: 4 x 4
  block trials  easy  hard
  <int>  <int> <dbl> <dbl>
1     1     25  1     0.8 
2     2     25  0.92  0.96
3     3     25  1     0.88
4     4     25  0.96  0.8 
\end{verbatim}

Inline prompt: ``Overall accuracy = NA. Easy = NA, Hard = NA.'' Replace
NA with inline code.

\begin{Shaded}
\begin{Highlighting}[]
\NormalTok{Overall accuracy }\OtherTok{=} \StringTok{\textquotesingle{}r round(mean(acc$acc),3)\textquotesingle{}}
\NormalTok{Easy }\OtherTok{=} \StringTok{\textquotesingle{}r round(mean(acc$acc[acc$condition == "easy"]), 3)\textquotesingle{}}
\NormalTok{Hard }\OtherTok{=} \StringTok{\textquotesingle{}r round(mean(acc$acc[acc$condition == "hard"]), 3)\textquotesingle{}}
\end{Highlighting}
\end{Shaded}

\begin{verbatim}
Error in parse(text = input): <text>:1:9: unexpected symbol
1: Overall accuracy
            ^
\end{verbatim}

\subsection{C. Reaction times (RT): shifted
lognormal}\label{c.-reaction-times-rt-shifted-lognormal}

RTs are famously right-skewed (ie., sometimes participants take a very
long time to respond). This looks really different than a normal
distribution. We'll use \texttt{t0\ +\ LogNormal(μ,\ σ)} to
approximately simulate their distribution. \texttt{t0} can be thought of
a ``baseline'' reaction time that could vary by participant.

\subsubsection{Simulating reaction
times}\label{simulating-reaction-times}

Tasks:

\begin{itemize}
\tightlist
\item
  Simulate \textbf{n = 500} RTs with \texttt{t0\ =\ 300},
  \texttt{meanlog\ =\ log(250)}, \texttt{sdlog\ =\ 0.35} into
  \texttt{rt\_df}.
\item
  Plot a histogram and compute \textbf{mean}, \textbf{median},
  \textbf{max}.
\item
  Create \texttt{log\_rt\ =\ log(rt)} and compare distributions
\end{itemize}

Here's a helper function to create lognormal RTs with a baseline shift.
Note on these parameters: In a lognormal distribution, the parameters
\texttt{meanlog} and \texttt{sdlog} describe the underlying normal
distribution on the log scale, NOT the mean and standard deviation of
the observed reaction times. The median of a lognormal is
\texttt{exp(meanlog)}, so \texttt{meanlog} = log(250) makes the median
of the lognormal part 250 ms. The helper adds \texttt{t0} on top, so
with \texttt{t0\ =\ 300} the median simulated RT is about 300 + 250 =
550 ms. The \texttt{sdlog} parameter controls the amount of skew on the
log scale: typical behavioral data are well-modeled by \texttt{sdlog}
values between about 0.3 and 0.5. Using \texttt{sdlog} = log(20) would
be incorrect, because \texttt{sdlog} already represents variability on
the log scale; taking the log of a raw-scale SD would produce
unrealistically long-tailed RTs spanning thousands of milliseconds. You
don't need to understand all of these details, but I realize it could be
confusing and so wanted to include this explanation.

\begin{Shaded}
\begin{Highlighting}[]
\CommentTok{\# Helper: shifted lognormal RTs (t0 + LogNormal(meanlog, sdlog))}
\NormalTok{r\_shifted\_lognorm }\OtherTok{\textless{}{-}} \ControlFlowTok{function}\NormalTok{(n, t0, meanlog, sdlog) \{}
\NormalTok{  t0 }\SpecialCharTok{+} \FunctionTok{rlnorm}\NormalTok{(n, }\AttributeTok{meanlog =}\NormalTok{ meanlog, }\AttributeTok{sdlog =}\NormalTok{ sdlog)}
\NormalTok{\}}
\end{Highlighting}
\end{Shaded}

\begin{Shaded}
\begin{Highlighting}[]
\CommentTok{\# }\AlertTok{TODO}\CommentTok{: Simulate RTs and their logged RTs}
\CommentTok{\# Hint: Use the r\_shifted\_lognorm() function with the specified parameters}
\CommentTok{\# Create a tibble with rt and log\_rt columns}

\CommentTok{\# Complete this structure:}
\NormalTok{rt\_df }\OtherTok{\textless{}{-}} \FunctionTok{tibble}\NormalTok{(}
  \AttributeTok{rt =} \FunctionTok{r\_shifted\_lognorm}\NormalTok{(}\DecValTok{500}\NormalTok{, }\AttributeTok{t0 =} \DecValTok{300}\NormalTok{, }\AttributeTok{meanlog =} \FunctionTok{log}\NormalTok{(}\DecValTok{250}\NormalTok{), }\AttributeTok{sdlog =}\NormalTok{ .}\DecValTok{35}\NormalTok{)  }\CommentTok{\# Fill in the parameters}
\NormalTok{) }\SpecialCharTok{|\textgreater{}}
  \FunctionTok{mutate}\NormalTok{(}\AttributeTok{log\_rt =} \FunctionTok{log}\NormalTok{(rt))  }\CommentTok{\# Fill in the column name}

\CommentTok{\# Check the first few rows}
\FunctionTok{head}\NormalTok{(rt\_df)}
\end{Highlighting}
\end{Shaded}

\begin{verbatim}
# A tibble: 6 x 2
     rt log_rt
  <dbl>  <dbl>
1  473.   6.16
2  430.   6.06
3  476.   6.17
4  568.   6.34
5  574.   6.35
6  521.   6.26
\end{verbatim}

Now look at their relationship -- use \texttt{ggplot()} with
\texttt{geom\_point()}, putting \texttt{rt} on the x-axis and
\texttt{log\_rt} on the y-axis.

\begin{Shaded}
\begin{Highlighting}[]
\CommentTok{\# }\AlertTok{TODO}\CommentTok{: Plot rt vs log\_rt}
\CommentTok{\# Hint: ggplot() with geom\_point()}
\FunctionTok{ggplot}\NormalTok{(rt\_df, }\FunctionTok{aes}\NormalTok{(}\AttributeTok{x =}\NormalTok{ rt, }\AttributeTok{y =}\NormalTok{ log\_rt)) }\SpecialCharTok{+}  \CommentTok{\# Fill in x and y column names}
  \FunctionTok{geom\_point}\NormalTok{(}\AttributeTok{alpha =} \FloatTok{0.5}\NormalTok{) }\SpecialCharTok{+}
  \FunctionTok{labs}\NormalTok{(}\AttributeTok{x =} \StringTok{"RT (ms)"}\NormalTok{, }\AttributeTok{y =} \StringTok{"Log RT"}\NormalTok{)}
\end{Highlighting}
\end{Shaded}

\pandocbounded{\includegraphics[keepaspectratio]{distributions-lab-intermediate_files/figure-pdf/unnamed-chunk-13-1.pdf}}

Now plot a histogram of \texttt{rt}, and another of \texttt{log\_rt}.
Which one looks more like a normal distribution?

\begin{Shaded}
\begin{Highlighting}[]
\CommentTok{\# Complete this structure:}
\FunctionTok{ggplot}\NormalTok{(rt\_df, }\FunctionTok{aes}\NormalTok{(}\AttributeTok{x =}\NormalTok{ rt)) }\SpecialCharTok{+}  \CommentTok{\# first rt, then copy the chunk and use log\_rt}
  \FunctionTok{geom\_histogram}\NormalTok{(}\AttributeTok{bins =} \DecValTok{40}\NormalTok{) }\SpecialCharTok{+}
  \FunctionTok{labs}\NormalTok{(}\AttributeTok{x =} \StringTok{"RT (ms)"}\NormalTok{, }\AttributeTok{y =} \StringTok{"Count"}\NormalTok{)}
\end{Highlighting}
\end{Shaded}

\pandocbounded{\includegraphics[keepaspectratio]{distributions-lab-intermediate_files/figure-pdf/unnamed-chunk-14-1.pdf}}

\begin{Shaded}
\begin{Highlighting}[]
\FunctionTok{ggplot}\NormalTok{(rt\_df, }\FunctionTok{aes}\NormalTok{(}\AttributeTok{x =}\NormalTok{ log\_rt)) }\SpecialCharTok{+}  \CommentTok{\# first rt, then copy the chunk and use log\_rt}
  \FunctionTok{geom\_histogram}\NormalTok{(}\AttributeTok{bins =} \DecValTok{40}\NormalTok{) }\SpecialCharTok{+}
  \FunctionTok{labs}\NormalTok{(}\AttributeTok{x =} \StringTok{"RT (ms)"}\NormalTok{, }\AttributeTok{y =} \StringTok{"Count"}\NormalTok{)}
\end{Highlighting}
\end{Shaded}

\pandocbounded{\includegraphics[keepaspectratio]{distributions-lab-intermediate_files/figure-pdf/unnamed-chunk-14-2.pdf}}

\begin{Shaded}
\begin{Highlighting}[]
\CommentTok{\# the log\_rt looks more closer to normal distribution.}
\end{Highlighting}
\end{Shaded}

OK, now let's make some summaries. Use the summarize function to
calculate the mean, median, max, and sd of the rt distribution. Note
that you don't have to make a new dataframe for this.

\begin{Shaded}
\begin{Highlighting}[]
\CommentTok{\# }\AlertTok{TODO}\CommentTok{: summaries}
\CommentTok{\# Hint: Use summarise() to calculate mean\_rt, median\_rt, max\_rt, sd\_rt}
\CommentTok{\# Complete this structure:}
\NormalTok{rt\_df\_summary }\OtherTok{\textless{}{-}}\NormalTok{ rt\_df }\SpecialCharTok{|\textgreater{}}
  \FunctionTok{summarise}\NormalTok{(}
    \AttributeTok{mean\_rt =} \FunctionTok{mean}\NormalTok{(rt),    }\CommentTok{\# Fill in column name}
    \AttributeTok{median\_rt =} \FunctionTok{median}\NormalTok{(rt), }\CommentTok{\# Fill in column name}
    \AttributeTok{max\_rt =} \FunctionTok{max}\NormalTok{(rt),      }\CommentTok{\# Fill in column name}
    \AttributeTok{sd\_rt =} \FunctionTok{sd}\NormalTok{(rt)         }\CommentTok{\# Fill in column name}
\NormalTok{  )}

\NormalTok{log\_rt\_df\_summary }\OtherTok{\textless{}{-}}\NormalTok{ rt\_df }\SpecialCharTok{|\textgreater{}}
  \FunctionTok{summarise}\NormalTok{(}
    \AttributeTok{mean\_rt =} \FunctionTok{mean}\NormalTok{(log\_rt),    }\CommentTok{\# Fill in column name}
    \AttributeTok{median\_rt =} \FunctionTok{median}\NormalTok{(log\_rt), }\CommentTok{\# Fill in column name}
    \AttributeTok{max\_rt =} \FunctionTok{max}\NormalTok{(log\_rt),      }\CommentTok{\# Fill in column name}
    \AttributeTok{sd\_rt =} \FunctionTok{sd}\NormalTok{(log\_rt)         }\CommentTok{\# Fill in column name}
\NormalTok{  )}

\CommentTok{\# Print the results}
\FunctionTok{print}\NormalTok{(rt\_df\_summary)}
\end{Highlighting}
\end{Shaded}

\begin{verbatim}
# A tibble: 1 x 4
  mean_rt median_rt max_rt sd_rt
    <dbl>     <dbl>  <dbl> <dbl>
1    564.      548.  1268.  101.
\end{verbatim}

\begin{Shaded}
\begin{Highlighting}[]
\FunctionTok{print}\NormalTok{(log\_rt\_df\_summary)}
\end{Highlighting}
\end{Shaded}

\begin{verbatim}
# A tibble: 1 x 4
  mean_rt median_rt max_rt sd_rt
    <dbl>     <dbl>  <dbl> <dbl>
1    6.32      6.31   7.15 0.167
\end{verbatim}

\subsection{D. Two conditions: valid vs invalid
cue}\label{d.-two-conditions-valid-vs-invalid-cue}

Backstory. In a spatial attention task (think Posner cueing), each trial
begins with a cue that is either valid (points to the upcoming target
location) or invalid (points away). Valid cues help orient attention and
typically speed responses; invalid cues slow responses. Reaction times
(RTs) are right-skewed, as in the prior example. Let's try to simulate
valid and invalid trials.

Add a \textbf{condition} factor. Suppose invalid cues slow responses by
\textasciitilde{}\textbf{40 ms} on average.

\subsubsection{Try it}\label{try-it}

\begin{Shaded}
\begin{Highlighting}[]
\CommentTok{\# }\AlertTok{TODO}\CommentTok{: Simulate RTs for valid and invalid conditions}
\CommentTok{\# Hint: Use r\_shifted\_lognorm() with different parameters for each condition}
\CommentTok{\# Invalid cues should be slower (higher meanlog)}

\CommentTok{\# Complete this structure:}
\NormalTok{rt\_df\_valid\_invalid }\OtherTok{\textless{}{-}} \FunctionTok{tibble}\NormalTok{(}
  \AttributeTok{rt\_valid =} \FunctionTok{r\_shifted\_lognorm}\NormalTok{(}\DecValTok{500}\NormalTok{, }\AttributeTok{t0 =} \DecValTok{300}\NormalTok{, }\AttributeTok{meanlog =} \FunctionTok{log}\NormalTok{(}\DecValTok{250}\NormalTok{), }\AttributeTok{sdlog =}\NormalTok{ .}\DecValTok{35}\NormalTok{),     }\CommentTok{\# Fill in parameters for valid condition}
  \AttributeTok{rt\_invalid =} \FunctionTok{r\_shifted\_lognorm}\NormalTok{(}\DecValTok{500}\NormalTok{, }\AttributeTok{t0 =} \DecValTok{300}\NormalTok{, }\AttributeTok{meanlog =} \FunctionTok{log}\NormalTok{(}\DecValTok{290}\NormalTok{), }\AttributeTok{sdlog =}\NormalTok{ .}\DecValTok{35}\NormalTok{)   }\CommentTok{\# Fill in parameters for invalid condition (slower)}
\NormalTok{)}

\CommentTok{\# Check the first few rows}
\FunctionTok{head}\NormalTok{(rt\_df\_valid\_invalid)}
\end{Highlighting}
\end{Shaded}

\begin{verbatim}
# A tibble: 6 x 2
  rt_valid rt_invalid
     <dbl>      <dbl>
1     623.       682.
2     559.       850.
3     578.       706.
4     471.       683.
5     549.       911.
6     453.       527.
\end{verbatim}

Use \texttt{pivot\_longer} to make this a long dataframe, where
\texttt{rt} is just one column, and you have \texttt{condition} in
another column. Use the \texttt{?} to look at the arguments for the
function.

\begin{Shaded}
\begin{Highlighting}[]
\CommentTok{\# }\AlertTok{TODO}\CommentTok{: Use pivot\_longer to reshape the data}
\CommentTok{\# Hint: Convert rt\_valid and rt\_invalid columns into rt and condition columns}
\CommentTok{\# Complete this structure:}
\NormalTok{rt\_df\_valid\_invalid\_longer }\OtherTok{\textless{}{-}}\NormalTok{ rt\_df\_valid\_invalid }\SpecialCharTok{|\textgreater{}}
  \FunctionTok{pivot\_longer}\NormalTok{(}
    \AttributeTok{cols =} \FunctionTok{c}\NormalTok{(}\StringTok{\textquotesingle{}rt\_valid\textquotesingle{}}\NormalTok{, }\StringTok{\textquotesingle{}rt\_invalid\textquotesingle{}}\NormalTok{),        }\CommentTok{\# Fill in the column names to pivot}
    \AttributeTok{values\_to =} \StringTok{\textquotesingle{}rt\textquotesingle{}}\NormalTok{,           }\CommentTok{\# Fill in the name for the values column}
    \AttributeTok{names\_to =} \StringTok{\textquotesingle{}condition\textquotesingle{}}             \CommentTok{\# Fill in the name for the names column}
\NormalTok{  )}

\CommentTok{\# Check the first few rows}
\FunctionTok{head}\NormalTok{(rt\_df\_valid\_invalid\_longer)}
\end{Highlighting}
\end{Shaded}

\begin{verbatim}
# A tibble: 6 x 2
  condition     rt
  <chr>      <dbl>
1 rt_valid    623.
2 rt_invalid  682.
3 rt_valid    559.
4 rt_invalid  850.
5 rt_valid    578.
6 rt_invalid  706.
\end{verbatim}

Now, use \texttt{group\_by(condition)} before computing your summaries
in tidyverse -- you should see outputs grouped by condition.

\begin{Shaded}
\begin{Highlighting}[]
\CommentTok{\# }\AlertTok{TODO}\CommentTok{: group\_by() condition and summarise mean\_rt, sd\_rt}
\CommentTok{\# Hint: Use group\_by(condition) then summarise()}
\CommentTok{\# Complete this structure:}
\NormalTok{valid\_invalid\_summary }\OtherTok{\textless{}{-}}\NormalTok{ rt\_df\_valid\_invalid\_longer }\SpecialCharTok{|\textgreater{}}
  \FunctionTok{group\_by}\NormalTok{(condition) }\SpecialCharTok{|\textgreater{}}  \CommentTok{\# Fill in the grouping variable}
  \FunctionTok{summarise}\NormalTok{(}
    \AttributeTok{mean\_rt =} \FunctionTok{mean}\NormalTok{(rt),  }\CommentTok{\# Fill in column name}
    \AttributeTok{sd\_rt =} \FunctionTok{sd}\NormalTok{(rt),      }\CommentTok{\# Fill in column name}
    \AttributeTok{n =} \FunctionTok{n}\NormalTok{()}
\NormalTok{  )}

\CommentTok{\# Print the results}
\FunctionTok{print}\NormalTok{(valid\_invalid\_summary)}
\end{Highlighting}
\end{Shaded}

\begin{verbatim}
# A tibble: 2 x 4
  condition  mean_rt sd_rt     n
  <chr>        <dbl> <dbl> <int>
1 rt_invalid    604. 109.    500
2 rt_valid      562.  96.6   500
\end{verbatim}

Now, make some sort of visual of the two distributions. Give the raw
dataframe to ggplot, and use \texttt{geom\_density} to plot the density
(you can toggle \texttt{alpha} here to change the transparency). Use
\texttt{fill=condition} within \texttt{aes} to specify that the two
conditions should have different colors.

\begin{Shaded}
\begin{Highlighting}[]
\DocumentationTok{\#\# }\AlertTok{TODO}\DocumentationTok{: Use ggplot to make density plot of simulated data}
\CommentTok{\# Hint: Use geom\_density() with fill=condition in aes()}
\CommentTok{\# Complete this structure:}
\FunctionTok{ggplot}\NormalTok{(rt\_df\_valid\_invalid\_longer, }\FunctionTok{aes}\NormalTok{(}\AttributeTok{x =}\NormalTok{ rt, }\AttributeTok{fill =}\NormalTok{ condition)) }\SpecialCharTok{+}  \CommentTok{\# Fill in dataframe and column names}
  \FunctionTok{geom\_density}\NormalTok{(}\AttributeTok{alpha =} \FloatTok{0.5}\NormalTok{) }\SpecialCharTok{+}            \CommentTok{\# Fill in alpha value (0{-}1)}
  \FunctionTok{labs}\NormalTok{(}\AttributeTok{title =} \StringTok{"Cueing Cost as Shift in Distribution"}\NormalTok{, }\AttributeTok{x =} \StringTok{"RT (ms)"}\NormalTok{, }\AttributeTok{y =} \StringTok{"Density"}\NormalTok{)}
\end{Highlighting}
\end{Shaded}

\pandocbounded{\includegraphics[keepaspectratio]{distributions-lab-intermediate_files/figure-pdf/unnamed-chunk-19-1.pdf}}

BONUS. Okay, so far you've been simulating ``raw'' data -- i.e., 500
trials of valid, and 500 trials of invalid cueing. Try simulating 20
different participants, where each participant completes 200 trials (100
valid, 100 invalid).

\begin{Shaded}
\begin{Highlighting}[]
\CommentTok{\# }\AlertTok{TODO}\CommentTok{: Simulate multiple participants}
\CommentTok{\# Hint: Use a for loop to iterate over participants}
\CommentTok{\# Each participant should have different baseline shifts}
\CommentTok{\# Use bind\_rows() to combine all participant data}

\CommentTok{\# Set up parameters}
\NormalTok{num\_participants }\OtherTok{\textless{}{-}} \DecValTok{20}
\NormalTok{num\_trials\_per\_condition }\OtherTok{\textless{}{-}} \DecValTok{100}

\CommentTok{\# Start with an empty tibble}
\NormalTok{df\_experiment }\OtherTok{\textless{}{-}} \FunctionTok{tibble}\NormalTok{()}

\CommentTok{\# Loop over participants}
\ControlFlowTok{for}\NormalTok{ (i }\ControlFlowTok{in} \DecValTok{1}\SpecialCharTok{:}\NormalTok{num\_participants) \{}
  \CommentTok{\# Random baseline shift for this participant }
\NormalTok{  baseline\_shift }\OtherTok{\textless{}{-}} \FunctionTok{rnorm}\NormalTok{(}\DecValTok{1}\NormalTok{, }\AttributeTok{mean =} \DecValTok{0}\NormalTok{, }\AttributeTok{sd =} \DecValTok{80}\NormalTok{)}
  
  \CommentTok{\# Simulate trials for this participant}
\NormalTok{  sim\_subj }\OtherTok{\textless{}{-}} \FunctionTok{tibble}\NormalTok{(}
    \AttributeTok{participant =}\NormalTok{ i,}
    \AttributeTok{condition =} \FunctionTok{rep}\NormalTok{(}\FunctionTok{c}\NormalTok{(}\StringTok{"valid"}\NormalTok{, }\StringTok{"invalid"}\NormalTok{), }\AttributeTok{each =}\NormalTok{ num\_trials\_per\_condition),}
    \AttributeTok{trial =} \FunctionTok{rep}\NormalTok{(}\DecValTok{1}\SpecialCharTok{:}\NormalTok{num\_trials\_per\_condition, }\AttributeTok{times =} \DecValTok{2}\NormalTok{),}
    \AttributeTok{rt =} \FunctionTok{c}\NormalTok{(}
      \FunctionTok{r\_shifted\_lognorm}\NormalTok{(num\_trials\_per\_condition, }\AttributeTok{t0 =} \DecValTok{300} \SpecialCharTok{+}\NormalTok{ baseline\_shift, }\AttributeTok{meanlog =} \FunctionTok{log}\NormalTok{(}\DecValTok{250}\NormalTok{), }\AttributeTok{sdlog =} \FloatTok{0.35}\NormalTok{),}
      \FunctionTok{r\_shifted\_lognorm}\NormalTok{(num\_trials\_per\_condition, }\AttributeTok{t0 =} \DecValTok{300} \SpecialCharTok{+}\NormalTok{ baseline\_shift, }\AttributeTok{meanlog =} \FunctionTok{log}\NormalTok{(}\DecValTok{290}\NormalTok{), }\AttributeTok{sdlog =} \FloatTok{0.35}\NormalTok{)}
\NormalTok{    )}
\NormalTok{  ) }\SpecialCharTok{|\textgreater{}}
    \FunctionTok{mutate}\NormalTok{(}\AttributeTok{log\_rt =} \FunctionTok{log}\NormalTok{(rt))}
  
  \CommentTok{\# Append to growing data frame}
\NormalTok{  df\_experiment }\OtherTok{\textless{}{-}} \FunctionTok{bind\_rows}\NormalTok{(df\_experiment, sim\_subj)}
\NormalTok{\}}

\CommentTok{\# Check the structure}
\FunctionTok{head}\NormalTok{(df\_experiment)}
\end{Highlighting}
\end{Shaded}

\begin{verbatim}
# A tibble: 6 x 5
  participant condition trial    rt log_rt
        <int> <chr>     <int> <dbl>  <dbl>
1           1 valid         1  355.   5.87
2           1 valid         2  409.   6.01
3           1 valid         3  380.   5.94
4           1 valid         4  350.   5.86
5           1 valid         5  453.   6.12
6           1 valid         6  420.   6.04
\end{verbatim}

\begin{Shaded}
\begin{Highlighting}[]
\CommentTok{\# }\AlertTok{TODO}\CommentTok{: Summarize by participant and condition}
\CommentTok{\# Hint: Use group\_by(participant, condition) and summarise()}
\CommentTok{\# Complete this structure:}
\NormalTok{df\_experiment\_summary }\OtherTok{\textless{}{-}}\NormalTok{ df\_experiment }\SpecialCharTok{|\textgreater{}}
  \FunctionTok{group\_by}\NormalTok{(participant, condition) }\SpecialCharTok{|\textgreater{}}  \CommentTok{\# Fill in grouping variables}
  \FunctionTok{summarise}\NormalTok{(}
    \AttributeTok{mean\_rt =} \FunctionTok{mean}\NormalTok{(rt),  }\CommentTok{\# Fill in column name}
    \AttributeTok{n\_trials =} \FunctionTok{n}\NormalTok{(),}
    \AttributeTok{.groups =} \StringTok{"drop"}
\NormalTok{  )}

\CommentTok{\# Print the results}
\FunctionTok{print}\NormalTok{(df\_experiment\_summary)}
\end{Highlighting}
\end{Shaded}

\begin{verbatim}
# A tibble: 40 x 4
   participant condition mean_rt n_trials
         <int> <chr>       <dbl>    <int>
 1           1 invalid      496.      100
 2           1 valid        441.      100
 3           2 invalid      680.      100
 4           2 valid        618.      100
 5           3 invalid      605.      100
 6           3 valid        568.      100
 7           4 invalid      725.      100
 8           4 valid        658.      100
 9           5 invalid      571.      100
10           5 valid        506.      100
# i 30 more rows
\end{verbatim}

BONUS: Plot average RTs and 95 CIs over individual subjects. You can
reuse \texttt{stat\_summary} arguments for the mean and CIs from above.
Plot raw data points from each participant using \texttt{geom\_point}.

\begin{Shaded}
\begin{Highlighting}[]
\CommentTok{\# }\AlertTok{TODO}\CommentTok{: Plot average RTs and 95 CIs over individual subjects}
\CommentTok{\# Hint: Use stat\_summary() for means and CIs, geom\_point() for individual data}
\CommentTok{\# Complete this structure:}
\FunctionTok{ggplot}\NormalTok{(df\_experiment\_summary, }\FunctionTok{aes}\NormalTok{(}\AttributeTok{x =}\NormalTok{ condition, }\AttributeTok{y =}\NormalTok{ mean\_rt, }\AttributeTok{color =}\NormalTok{ condition, }\AttributeTok{fill =}\NormalTok{ condition)) }\SpecialCharTok{+}  \CommentTok{\# Fill in dataframe and aesthetics}
  \FunctionTok{geom\_point}\NormalTok{(}\AttributeTok{alpha =} \FloatTok{0.6}\NormalTok{) }\SpecialCharTok{+}  \CommentTok{\# Individual subject means}
  \FunctionTok{geom\_line}\NormalTok{(}\FunctionTok{aes}\NormalTok{(}\AttributeTok{group =}\NormalTok{ participant), }\AttributeTok{alpha =} \FloatTok{0.3}\NormalTok{, }\AttributeTok{color =} \StringTok{\textquotesingle{}grey\textquotesingle{}}\NormalTok{) }\SpecialCharTok{+}  \CommentTok{\# Connect data from each participant}
  \FunctionTok{stat\_summary}\NormalTok{(}\AttributeTok{fun =}\NormalTok{ mean, }\AttributeTok{geom =} \StringTok{"point"}\NormalTok{, }\AttributeTok{size =} \DecValTok{3}\NormalTok{) }\SpecialCharTok{+}  \CommentTok{\# Overall mean}
  \FunctionTok{stat\_summary}\NormalTok{(}\AttributeTok{fun.data =}\NormalTok{ mean\_se, }\AttributeTok{fun.args =} \FunctionTok{list}\NormalTok{(}\AttributeTok{mult =} \FloatTok{1.96}\NormalTok{), }\AttributeTok{geom =} \StringTok{"errorbar"}\NormalTok{, }\AttributeTok{width =} \FloatTok{0.1}\NormalTok{) }\SpecialCharTok{+}  \CommentTok{\# 95\% CIs}
  \FunctionTok{labs}\NormalTok{(}\AttributeTok{x =} \StringTok{"Condition"}\NormalTok{, }\AttributeTok{y =} \StringTok{"Average RT"}\NormalTok{, }\AttributeTok{title =} \StringTok{"RT by Simulated Cueing Condition"}\NormalTok{)}
\end{Highlighting}
\end{Shaded}

\pandocbounded{\includegraphics[keepaspectratio]{distributions-lab-intermediate_files/figure-pdf/unnamed-chunk-22-1.pdf}}




\end{document}
